\runninghead{The Lights in the Church}
\properlane{coming-through-the-church}{The Lights of Heaven in the Church: Christ Comes Through the Age of Trial, and the Hour Is Hidden}

The Gospel names three lights, \latin{sol, luna, stellæ}, and says only that
there will be signs in them. St Matthew's account of the same discourse is
more specific: ``the sun shall be darkened and the moon shall not give her
light and the stars shall fall from heaven'' (Mt 24:29). In St Augustine's
time a bishop, Hesychius, wrote to him that these words were being fulfilled
before their eyes: no country and no place was free of affliction, so men
were plainly ``withering away for fear'' (Letter 199.36). Augustine's reply
turns on what the sun, the moon and the stars are. He and St Ambrose, who
expounded the same verses of Luke, both read the lights as figures of the
Church under persecution, and both find Christ's coming already under way
within the Church's history. They assign the sun, the moon and the stars
differently.

\subsection{Ambrose: the sun in the believer's faith}

Ambrose expounds Luke's verse with Matthew's fuller wording before him
(\work{Expositio evangelii secundum Lucam} X.37). ``When very many fall away
from religion, bright faith will be darkened by the cloud of
faithlessness, \latin{quia mihi sol ille coelestis mea fide vel minuitur, vel
augetur}, because that heavenly sun is lessened or increased for me by my
faith.'' He explains by the visible sun. Many look at its rays, and it seems
paler or brighter according to each one's power of sight; ``so spiritual
light is poured into each according to the devotion of the believer.'' The
darkened sun is Christ, whose light is dimmed in those who abandon the faith.

The moon is the Church. In its monthly course the moon fades when the earth
comes between it and the sun; ``so too the holy Church, when the vices of the
flesh stand in the way of the heavenly light, cannot borrow the brightness of
divine light from the rays of Christ.'' Ambrose names the vice that does this
in time of persecution: \latin{in persecutionibus lucem Dei solus plerumque
amor vitae huius excludit}, ``in persecutions it is often love of this life
alone that shuts out the light of God'' (X.37).

The stars that fall are \latin{illi utique resurrectionis gloria iam
micantes}, ``those already shining with the glory of the resurrection,'' the
men of whom God said to Abraham that his seed would be as the stars. ``So the
patriarchs will fall from men's sight, the prophets will fall, if the
bitterness of persecution grows strong.'' These things must happen ``until
the fullness of the virtues is heaped up in the Church, in all and in each,''
and Ambrose gives the reason: \latin{sic enim probantur boni, sic produntur
infirmi}, ``for so the good are proved, and so the weak are exposed''
(X.38). The withering of men in the next verse he takes of the soul: under
fear of judgment ``the dew of the sacred fountain dries up in us; for
faithlessness is dry, and faith is dewy'' (X.38).

When he reaches the coming of the Son of man, Ambrose offers a reading with
the word \latin{Fortasse}, ``Perhaps.'' As the Lord's coming is awaited, so
that his presence may be in the whole world, \latin{quae fit utique in
singulis, cum omnibus affectibus receperint Christum}, ``a presence that
surely comes about in each one, when they have received Christ with all their
affections,'' so the powers of heaven are moved by an increase of grace
(X.39). He refuses to imagine Christ coming in a rain cloud: ``I do not think
that Christ comes in murky darkness and rainy dread \dots\ How did he set his
tabernacle in the sun, if his coming will be rainy?'' (X.41). The clouds are
the prophets and saints, through whom Christ came, and Christ came \latin{in
nube levi incarnatus ex Virgine}, ``in a light cloud, made flesh of the
Virgin,'' for ``the prophet saw as it were a cloud coming from the East'' (Isa
19:1). The cloud was light, Ambrose says, because ``earthly vices did not weigh it down''
(X.42).

\subsection{Augustine: the Church is sun, moon and stars}

In his letter to Hesychius Augustine first sorts the discourse. Asked about
the fall of the Temple and the sign of his coming, Christ foretold what would
follow from that time, and the signs he gave concern three things: the fall
of the city, his coming in his body the Church, and his coming at the end in
the Church's head, himself (Letter 199.25--26). In the Church, Augustine
says, he does not cease to come until the end, and he is recognised coming in
his own \latin{dum eius quotidie membra nascuntur}, ``while his members are
born daily'' (199.25). Luke, he observes, keeps the siege of the city apart
from the last coming (199.27).

He then weighs Hesychius's claim against what history records. Were the signs
in heaven and earth that their own generation had seen greater than those
before it? The histories of the nations report wonders so many that some are
not even believed, and writers recorded that two suns were seen before the
Lord came in the flesh. When, Augustine asks, has anyone seen the sun darkened
as it was ``when the Light of the world hung on the tree''? That darkness came
at the Passover, at the full moon, when by the astronomers' own reckoning the
sun cannot be eclipsed (199.34). As for wars, ``when has the earth not been
worn down by them?'' Under the emperor Gallienus, when barbarians overran the
Roman provinces, many Christians could have believed that the end was near,
long after the Ascension (199.35).

Nor do the present evils fit the Gospel's words. St Paul says that sudden
destruction will come upon men ``when they shall say, Peace and security''
(1~Thess 5:3). The tribulation the Gospel foretells, Augustine reasons, must
then fall on the Church, troubled throughout the world, and not on those who
trouble her, for they will be saying ``peace.'' But the evils of his own day,
he says, are common \latin{utrique genti et utrique regno}, ``to both peoples
and both kingdoms,'' Christ's and the devil's, and the good and the wicked
suffer them alike. Amid them the banquets, the drunkenness, the lyres and the
dice go on: \latin{hoccine est arescere prae timore, an potius madescere prae
libidine?} ``Is this withering for fear, or rather soaking in lust?''
(199.36--37).

From both arguments he draws his conclusion. The signs, he thinks, are better
understood of the Church, ``lest the Lord Jesus should seem to have foretold,
as a great thing for the approach of his second coming, what used to befall
this world even before his first coming, and we be laughed at by those who
have read of more such things, and far greater, in the histories of the
nations'' (199.39). Scripture then gives him the figure. He reads the three
lights together: \latin{Ecclesia est enim sol et luna et stellae}, ``for the
Church is sun and moon and stars,'' because the Song of Songs says of her,
\latin{Speciosa sicut luna, electa sicut sol}, ``fair as the moon, bright as
the sun'' (Cant 6:9). When persecutors rage beyond measure and the wicked say
``Peace and security,'' sun, moon and stars fail together, and
\latin{Ecclesia non apparebit}, ``the Church will not appear.'' The stars that
fall are \latin{multi, qui gratia fulgere videbantur}, ``many who seemed to
shine with grace,'' who will give way to the persecutors and fall (199.39).
The distress of nations is that of the nations who will stand on the left hand
at the judgment (199.40).

The cloud in which the Son of man comes may be the Church, in which he does
not cease to come even now, or his own body, as the angels said at the
Ascension that he would come as he had gone (Acts 1:11). Which of the two to
choose, Augustine says, is hard to judge (199.41--42). The plainer sense is the
coming in person; but Scripture must be searched, and he reads on to the fig
tree: ``when you shall see these things come to pass, know that the kingdom
of God is at hand.'' Among ``these things'' is the Son of man seen coming in
the cloud, so when he is seen, the kingdom is not yet come, only near
(199.42). Perhaps, then, the whole passage can be taken of the coming by which
\latin{quotidie venit in corpore suo, quod est Ecclesia}, ``he comes daily in
his body, which is the Church,'' coming in his members \latin{tamquam in nube
magna fructificante atque crescente in universo mundo}, ``as in a great
cloud, bearing fruit and growing in the whole world,'' except what plainly
belongs to the last coming (199.45).

\subsection{What the two Fathers share, and where they differ}

Ambrose's sun is Christ, received more or less brightly by each believer.
Augustine's sun is the Church, which is sun, moon and stars at once.
Ambrose's falling stars are the patriarchs and prophets, already glorious,
lost to men's sight when persecution is fierce. Augustine's are the many in
the Church who seemed to shine and then yielded.

They agree on the main thing. Both read the signs of the Church in
persecution, when many fall from the faith: \latin{plurimis \dots\ a
religione deficientibus}, Ambrose writes; \latin{multi \dots\ cedent et
cadent}, Augustine writes. In both, the Church's light is lost to sight in
that time, dimmed by the flesh for Ambrose, hidden for Augustine. And both let
the coming the signs announce be one already happening: in each person who
receives Christ with all his affections (Ambrose), and daily in his body the
Church (Augustine).

\subsection{This generation}

The lesson ends with a saying sealed by \latin{Amen}: ``this generation shall
not pass away till all things be fulfilled'' (Lk 21:32). St John Chrysostom,
preaching on the same saying in Matthew, puts the difficulty himself: ``How
then, one may ask, did He say, This generation?'' He reads the discourse of
real trials, the wars, famines, false Christs and the rest, ``which we said
were to occur until His coming,'' and he takes the generation as one that
outlasts them: ``Speaking not of the generation then living, but of that of
the believers. For He is wont to distinguish a generation not by times only,
but also by the mode of religious service,'' as the psalm says, ``This is the
generation of them that seek the Lord'' (Ps 23:6). Christ's promise is that
``the generation of the faithful shall remain, cut off by none of the things
that have been mentioned \dots\ over this generation shall nothing prevail,
not famine, not pestilence, not earthquake, nor the tumults of wars, not false
Christs, not false prophets'' (\work{Homilies on Matthew} 77). Chrysostom does
not read the lights as figures of the Church. What he says of the generation, a
people of believers that no trial cuts off, is said of the same Church that
Ambrose and Augustine see tried in every age. St Bede gives two other senses for \latin{generatio}: the
whole human race, or the Jewish people in particular (\work{In Lucae
evangelium expositio} VI, on 21:32).

\subsection{The hour is hidden}

Augustine draws the moral consequence in the opening of the same letter. No
one knows the time of the world's end, but each knows that his own will come:
\latin{in quo enim quemque invenerit suus novissimus dies, in hoc eum
comprehendet mundi novissimus dies}, ``in whatever state his own last day
finds each one, in that state the last day of the world will overtake him''
(199.2). The watch Christ commands just after the lesson (Lk 21:36) is
therefore daily. St Cyril of
Jerusalem gives the same counsel to his catechumens: ``as to the time let no
one be curious \dots\ venture not thou to declare when these things shall
be, nor on the other hand supinely slumber'' (\work{Catechetical Lectures}
15.4).

\subsection{The difficulty: figure and event}

This interpretation stands against the first. Gregory and the two Cyrils read
the signs as coming portents in the heavens; Ambrose and Augustine read them
chiefly as figures of the Church. The disagreement is real, and neither Father
of the second group removes the last coming. Ambrose's \latin{Fortasse} marks
his reading of the coming as tentative (X.39), and Augustine twice hesitates
and excepts what plainly concerns the end (199.42, 45). Cyril of Jerusalem's warning that
converts from the Manichees must not think ``that this sun which shall be
darkened is Christ'' (15.3) is aimed at worship of the visible sun as a god,
and it does not touch Ambrose, whose sun is Christ's light in the believer.

\subsection{The Church's psalm and her prayer for his coming}

Augustine hears Psalm 24 as the voice of the Church (\work{En.\ in Ps.}
24.1). If the Gospel's signs describe the Church's age of trial, the psalm the
Mass sings at the entrance, after the Epistle and again at the offering is her
voice in that age: she lifts her soul while enemies wait to laugh, and she
trusts that none who wait for God are confounded.

The Introit verse and the Gradual ask God to show and make known his ways.
St Thomas Aquinas explains the ways as the commandments and the paths as the
counsels, and adds of the ways, \latin{Per has enim vias venit dominus ad
nos}, ``for by these ways the Lord comes to us'' (\work{Super Psalmos} 24,
n.~4). That is a coming in the present, by the commandments and above all the
commandment of charity, and it is the
kind of coming Ambrose and Augustine find already under way in the Gospel.
The Collect follows the Introit, and its \latin{veni}, addressed to the Son,
may be prayed of that daily coming as well as of the last.

\subsection{The day that has come and still draws near}

St Paul's ``The night is passed, and the day is at hand'' (Rom 13:12) has
several senses for Aquinas. The night already past for those Paul wrote to
was, ``more truly,'' the night of sin, and the time of the Law before Christ
had passed as well. The day at hand Paul meant as the day of glory; it can
also be the time of Christ's grace, \latin{quod etsi iam advenerit secundum
temporum cursum, appropinquare tamen nobis dicitur per fidem et
devotionem}, ``which, though it has already come in the course of time, is
said to draw near to us by faith and devotion'' (\work{Super Epistolam ad
Romanos} 13, lect.~3). His day of grace and the coming that Ambrose and
Augustine find in the Gospel have the same shape: Christ has come, and he
still comes to each one who believes.

The Alleluia then asks, \latin{Osténde nobis, Dómine, misericórdiam tuam},
``Shew us, O Lord, thy mercy.'' Preaching on that verse, Augustine describes
those who sing it as sons of God already, \latin{sed per fidem: sumus quidem
in spe, nondum sumus in re}, ``but by faith; we are so in hope, not yet in
reality'' (\work{En.\ in Ps.} 84.9), and he quotes St John: \latin{nondum
apparuit quid erimus}, ``it hath not yet appeared what we shall be'' (1~Jn
3:2).

Ambrose's moon loses the sun's light when the vices of the flesh stand in the
way. Over the gifts the Secret asks that the holy gifts, cleansing us by their
power, bring us purer \latin{ad suum princípium}, to their source. If
Ambrose's figure is brought to this prayer, that cleansing is what would take
the obstacle away.

\subsection{The cloud that bears fruit}

Augustine's great cloud, in which Christ comes, bears fruit and grows in the
whole world. Earlier in the same letter he takes the clouds of the coming to
be those of which the prophet says, \latin{Mandabo nubibus meis ne pluant
super eam}, ``I will command my clouds not to rain upon it'' (Isa 5:6; Letter
199.25): clouds that carry rain, which God can command to withhold it. The
Communion sings, \latin{Dóminus dabit benignitátem: et terra nostra dabit
fructum suum}, ``The Lord will give goodness: and our earth shall
yield her fruit.'' Preaching on that verse, Augustine prays for God's rain on
the seed of the word in his hearers, so that the earth broken by confession
may yield good works (\work{En.\ in Ps.} 84.15). If the rain he prays for in
the sermon falls from the clouds of his coming in the letter, the harvest the
Communion sings is the fruit of Christ's coming in his Church, the cloud that
rains on the earth and makes it fruitful.

The Postcommunion asks to receive God's mercy \latin{in médio templi tui}, ``in
the midst of thy temple.'' Christ comes daily in his body, Augustine says;
the temple where the Church asks to receive his mercy would be one place of
that coming. The prayer still looks ahead, to \latin{ventúra solémnia}, the
solemnities to come.

\begin{fourSenses}
\item[Literal.] Christ foretells signs in sun, moon and stars and his coming
in a cloud, and promises that his words outlast heaven and earth and that
``this generation'' will not pass until all is fulfilled (Lk 21:25--33).
\item[Allegorical.] The signs figure the Church under persecution, in which
many fall from the faith. For Ambrose the sun is Christ as faith receives him,
the moon the Church, which the vices of the flesh eclipse, and the falling
stars the patriarchs and prophets, lost to men's sight; the cloud is the
prophets and saints, and the flesh taken of the Virgin (X.37--42). For
Augustine the Church is sun, moon and stars, and the falling stars are those
who seemed to shine and gave way; the cloud is the Church or Christ's own body
(Letter 199.39--41).
\item[Moral.] Keep faith, so that the sun is not darkened for you
(Ambrose); watch daily, because each one's last day overtakes him as he is
(Augustine); do not be curious about the time, and do not slumber (Cyril of
Jerusalem).
\item[Anagogical.] The coming in the head completes the coming in the body
(Augustine), and the generation of the faithful stands at the end, unharmed
by everything that has passed (Chrysostom).
\end{fourSenses}
